% Self-authored one-page probe for the textbook fixture's footnote placement.
\documentclass[10pt,oneside,openany]{book}
\usepackage{amsmath}
% Deliberately expose the format-profile difference: upLaTeX moves this
% tiny tail glue across the footnote, while standard LaTeX and PraTeX do not.
\raggedbottom
\def~{\leavevmode\nobreak\ }
\tracingoutput=1
\showboxbreadth=10000
\showboxdepth=10000
\begin{document}
\mainmatter
\chapter{Measured Transformations}
\section{Model and vocabulary}
A measured transformation begins with a state, accepts a controlled input,
and produces a new state. The purpose of the model is not to predict every
detail. It is to expose which quantities are preserved, which quantities are
supplied, and which observations would disprove the explanation. This module
uses a deliberately small vocabulary so that each later argument can be
checked from the same definitions.\footnote{The prose and numerical examples
in this fixture are synthetic and have no external source.}
\begin{equation}
  s_{k+1}=s_k+u_k-e_k .
\end{equation}
Here $s_k$ is the stored state, $u_k$ is a supplied increment, and $e_k$ is
an observed loss. Equation (1.1) is a bookkeeping identity; interpretation
enters only when the three terms are measured.
\end{document}
